\section{Discussion}
\label{sec:conclusion}

The three studies connect loss weighting, gradient concentration in a
shared MOPD student, and vocabulary-level supervision. Their interpretation
depends on distinguishing raw gradients from optimizer steps and
cumulative parameter changes. Sparse cumulative changes alone do not
establish sparse gradients, while overlapping high-magnitude gradients
can indicate either shared or opposing corrections.

The PG/top-$k$ comparison tests for a material difference rather than
assuming an advantage for sampled supervision. A resolved absence of
such a difference is informative within the tested regime; a stable
gap can motivate further analysis. Likewise, teacher distribution
differences may or may not predict which parameters receive strong
gradients. Capability measurements determine whether these optimization
patterns accompany useful integration.

\missing{State the supported observations, including null relationships,
and their model, task, teacher, loss, and measurement scope after
completing the experiments.}
